\section{Conclusion}
Cross-node Direct Dispatch and Combine coordinate their participants at both transfer
boundaries: peer readiness gates injection, and sender completion followed by a local
join precedes an EP-wide tail exchange. On two H200 nodes, each boundary costs
11--15\us{} at every batch size we measured, which is roughly 20\% and 18\% of the
operator at eight tokens per rank and falls monotonically to about 5\% by 128; in
Qwen3-30B-A3B decode, crossing the node boundary doubles the entry share at fixed EP8. The cost is fixed, and small-batch
decode is where it is exposed.

\sys{} addresses both boundaries by changing where coordination is resolved rather than
what it guarantees. Readiness aggregated at the fan-in point lets a rank upload as soon
as it is locally ready, with bounded staging absorbing early arrivals; a sealed
per-destination traffic plan lets the network append an ordered completion notification
behind that destination's final data unit, so the destination need not wait for a sender
to traverse its own completion path and originate a fresh exchange. Data visibility,
local completion conditions, and buffer ownership are preserved throughout; only the
placement of the completion decision changes.

What this study does and does not settle is worth stating plainly. The baselines are
measured and the coordination stages are established from source. The latency
opportunity is analytical: under a symmetric reference each eligible boundary is worth
approximately one one-way traversal, and we characterize how much added overhead erases
it. What remains open is the exposed fraction of the tail, which requires timestamping
destination visibility separately from sender completion, and the achievable saving,
which requires hardware. Building that prototype, and measuring $\epsilon$ against the
conditions specified here, is the natural next step.
